% Book 3, Year 1, chapter 14 -- dynamic arrays and amortised analysis.
% Provenance (never printed): sources/book3/BRIEFS.md, chapter 14.
% ledger: b1:cpp-vector-amortised
\chapter{Dynamic Arrays and Amortised Analysis}\label{ch:b1:dynamic-arrays}

A Python list grows by one element at a time, millions of times, and each \code{append} takes about the same short
time however long the list already is. Yet a list lives in one block of memory, and a block cannot grow in place: when
it is full, a larger block must be reserved and every element copied into it, which costs time proportional to the
length. Both statements are true. The way out is to make the copies rare, and to count the cost of a long sequence of
operations rather than of one: this is amortised analysis, and this chapter builds the growing \omterm{def:b1:pointers-and-arrays:array}{array} on which it
works, in C, behind an interface that hides how it is made.

\begin{recall}
A C module is a header, which declares what its users may call, and a \omterm{def:b1:compilation-pipeline:source-file}{source file}, which defines it; the users include
the header only (\cref{ch:b1:c-modules-make-tests}). \code{malloc} reserves a block on the heap and \code{free} returns
it (\cref{ch:b1:pointers-and-arrays}). A cost in $\Theta(1)$ is bounded by a constant independent of the \omterm{def:b1:asymptotic-analysis:cost}{input
size} (\cref{ch:b1:asymptotic-analysis}).
\end{recall}

\section{Abstract types and implementations}\label{sec:b1:dynamic-arrays:adt}

\begin{definition}[Abstract data type]\label{def:b1:dynamic-arrays:adt}
An \emph{abstract data type}\index{abstract data type} is a type given by a set of operations and their specification,
without saying how its values are represented. Its \emph{interface} is the list of the operations with their
signatures; an \emph{implementation} chooses a representation and writes the operations. Code that uses only the
interface works unchanged with any implementation that meets the specification.
\end{definition}

\begin{definition}[Constructor, accessor, transformer]\label{def:b1:dynamic-arrays:operations}
The operations of an abstract type are of three kinds: a \emph{constructor}\index{constructor (of an abstract type)}
creates a value of the type, an \emph{accessor}\index{accessor} returns information about a value without changing it,
and a \emph{transformer}\index{transformer} changes it.
\end{definition}

\omcode{code/bachelor-1/14-dynamic-arrays/listings/vec.h}{1}{17}{The interface of the \omterm{def:b1:dynamic-arrays:dynamic-array}{dynamic
array}}\label{lst:b1:dynamic-arrays:vec-h}

The declaration \code{typedef struct vec vec;} names a structure without giving its fields: a user can hold a
\code{vec *} and pass it to the functions, but cannot read or change the fields, which are defined only in the
implementation. The \omterm{def:b1:compilation-pipeline:compiler}{compiler} enforces the interface. Two files implement it: \code{vec.c}, the growing \omterm{def:b1:pointers-and-arrays:array}{array} of this
chapter, and \code{fixed.c}, an \omterm{def:b1:pointers-and-arrays:array}{array} of 16 cells that refuses a seventeenth element.

\begin{omfigure}
\begin{tikzpicture}[b/.style={draw, rounded corners, align=center, font=\small, minimum height=9mm}]
\node[b, fill=omCode!8, minimum width=3.6cm] (user) at (0,1.8) {\code{vec_demo.c}\\ \scriptsize includes \code{vec.h}
only};
\node[b, fill=omDef!15, minimum width=10.6cm, font=\scriptsize] (h) at (0,0.3) {\code{vec.h}: \code{vec_new}, \code{vec_size},
\code{vec_get}, \code{vec_set}, \code{vec_push}, \code{vec_pop}};
\node[b, fill=black!4, minimum width=4cm] (v) at (-2.6,-1.4) {\code{vec.c}\\ \scriptsize grows by doubling};
\node[b, fill=black!4, minimum width=4cm] (f) at (2.6,-1.4) {\code{fixed.c}\\ \scriptsize 16 cells, then an error};
\draw[-{Stealth}] (user) -- (h); \draw[-{Stealth}] (v) -- (h.south -| v); \draw[-{Stealth}] (f) -- (h.south -| f);
\node[font=\scriptsize, anchor=west] at (5.6,1.8) {the user};
\node[font=\scriptsize, anchor=west] at (5.6,0.3) {the interface};
\node[font=\scriptsize, anchor=west] at (5.6,-1.4) {two implementations};
\end{tikzpicture}
\par\smallskip{\small An abstract type: the user depends on the interface only, and either implementation can be
linked behind it. The fields of \code{struct vec} differ between the two files.}
\end{omfigure}

\omoutput{code/bachelor-1/14-dynamic-arrays/outputs/swap_impl.txt}{The same demonstration linked with each
implementation: the fixed one stops at its \omterm{def:b1:dynamic-arrays:dynamic-array}{capacity}.}\label{out:b1:dynamic-arrays:swap}

\section{The doubling array}\label{sec:b1:dynamic-arrays:doubling}

\begin{definition}[Dynamic array, size and capacity]\label{def:b1:dynamic-arrays:dynamic-array}
A \emph{dynamic array}\index{dynamic array} stores a sequence in a block of memory larger than needed. Its
\emph{size}\index{size (of a dynamic array)} is the number of elements, its \emph{capacity}\index{capacity (of a dynamic
array)} the number of cells in the block. Adding an element to a full \omterm{def:b1:pointers-and-arrays:array}{array} first moves the elements to a larger block,
whose capacity is a constant factor larger.
\end{definition}

\begin{definition}[Representation invariant]\label{def:b1:dynamic-arrays:rep-invariant}
A \emph{representation invariant}\index{representation invariant} is a condition on the fields of an implementation
that every operation may assume on entry and must restore before returning. It says which values of the fields
represent a value of the abstract type.
\end{definition}

\omcodes{code/bachelor-1/14-dynamic-arrays/listings/vec.c}{6-25,49-58}{The doubling \omterm{def:b1:pointers-and-arrays:array}{array}: the representation, the
resize, and the two operations that change the \omterm{def:b1:dynamic-arrays:dynamic-array}{size}}\label{lst:b1:dynamic-arrays:vec}

The \omterm{def:b1:dynamic-arrays:dynamic-array}{capacity} doubles when a push finds the \omterm{def:b1:pointers-and-arrays:array}{array} full, and halves when a pop leaves it a quarter full. Shrinking at a
half would be a mistake (\cref{sec:b1:dynamic-arrays:shrinking}). The function \code{resize} copies every element and
counts the copies, so that the experiments can report them.

\omoutput[1-5,8-10,16-24]{code/bachelor-1/14-dynamic-arrays/outputs/vec_demo.txt}{Seventeen pushes, then the pops
that shrink the \omterm{def:b1:pointers-and-arrays:array}{array}.}\label{out:b1:dynamic-arrays:demo}

\section{Amortised cost}\label{sec:b1:dynamic-arrays:amortised}

A push usually costs one write; the push that finds the \omterm{def:b1:pointers-and-arrays:array}{array} full costs one write plus as many copies as there are
elements. The worst case of one push is therefore $\Theta(n)$, which is true and misleading: such pushes are rare.

\begin{definition}[Amortised cost]\label{def:b1:dynamic-arrays:amortised}
Operations on a data structure have \emph{amortised cost}\index{amortised cost} $a$ if every sequence of $m$
operations, starting from the empty structure, costs at most $a \cdot m$ in total. The bound concerns the worst sequence,
not an average over inputs: no probability is involved.
\end{definition}

\begin{definition}[Aggregate, accounting and potential methods]\label{def:b1:dynamic-arrays:methods}
Three methods prove amortised bounds. The \emph{aggregate method}\index{aggregate method} bounds the total cost of any
sequence directly. The \emph{accounting method}\index{accounting method} charges each operation a fixed number of
credits, spends one credit per unit of actual cost, and shows that the credits stored never become negative. The
\emph{potential method} chooses a \emph{potential function}\index{potential function} $\Phi$ from the states of the
structure to the reals, and calls $c_i + \Phi_i - \Phi_{i-1}$ the \omterm{def:b1:dynamic-arrays:amortised}{amortised cost} of the $i$-th operation, of actual
cost $c_i$; if $\Phi_m \geq \Phi_0$ after every sequence, the total actual cost is at most the total \omterm{def:b1:dynamic-arrays:amortised}{amortised cost}.
\end{definition}

The last claim is a telescoping sum: $\sum_i (c_i + \Phi_i - \Phi_{i-1}) = \sum_i c_i + \Phi_m - \Phi_0 \geq \sum_i
c_i$.

\begin{proposition}[Pushes on the doubling array]\label{prop:b1:dynamic-arrays:pushes}
Starting from an empty \omterm{def:b1:pointers-and-arrays:array}{array} of \omterm{def:b1:dynamic-arrays:dynamic-array}{capacity} 1, $n$ pushes make fewer than $2n$ copies, and cost fewer than $3n$ writes and
copies in all: a push has \omterm{def:b1:dynamic-arrays:amortised}{amortised cost} 3.
\end{proposition}

\begin{proof}
\emph{Aggregate}: the copies happen when the \omterm{def:b1:dynamic-arrays:dynamic-array}{size} reaches $1, 2, 4, \dots, 2^j$ with $2^j < n$, and copy that many
elements: $1 + 2 + \dots + 2^j = 2^{j+1} - 1 < 2n$. With the $n$ writes, the total is below $3n$.

\emph{Accounting}: charge 3 credits per push. One pays its write; two are stored with the new element. When an \omterm{def:b1:pointers-and-arrays:array}{array}
of \omterm{def:b1:dynamic-arrays:dynamic-array}{capacity} $c$ fills up, the $c/2$ elements pushed since the last resize have stored $c$ credits, which pay for the $c$
copies. The store never goes negative.

\emph{Potential}: let $\Phi = 2 \cdot \text{size} - \text{capacity}$, which is $-1$ for the empty \omterm{def:b1:pointers-and-arrays:array}{array} and at least 0
after any push. A push without resize costs 1 and raises $\Phi$ by 2: amortised 3. A push on a full \omterm{def:b1:pointers-and-arrays:array}{array} of \omterm{def:b1:dynamic-arrays:dynamic-array}{size} $s$
costs $1 + s$, and $\Phi$ goes from $2s - s = s$ to $2(s + 1) - 2s = 2$: amortised $1 + s + 2 - s = 3$.
\end{proof}

\begin{omfigure}
\begin{tikzpicture}
\begin{axis}[width=10cm, height=5.6cm, xmin=0.3, xmax=17.7, ymin=0, ymax=34, xtick={1,2,3,5,9,17},
  xlabel={push number}, legend pos=outer north east, legend style={font=\small}, tick label style={font=\small}]
\addplot[const plot, thick, omDef] table[x=i, y=capacity]
  {code/bachelor-1/14-dynamic-arrays/figures/copies.dat}; \addlegendentry{capacity}
\addplot[thick, omCode, mark=*, mark size=1.2pt] table[x=i, y=size]
  {code/bachelor-1/14-dynamic-arrays/figures/copies.dat}; \addlegendentry{size}
\addplot[ycomb, very thick, omThm] table[x=i, y=cost]
  {code/bachelor-1/14-dynamic-arrays/figures/copies.dat}; \addlegendentry{actual cost of the push}
\addplot[thick, dashed] table[x=i, y=amortised]
  {code/bachelor-1/14-dynamic-arrays/figures/copies.dat}; \addlegendentry{amortised cost: 3}
\end{axis}
\end{tikzpicture}
\par\smallskip{\small Seventeen pushes on the doubling \omterm{def:b1:pointers-and-arrays:array}{array}, from \cref{out:b1:dynamic-arrays:demo}. The \omterm{def:b1:dynamic-arrays:dynamic-array}{capacity}
doubles at pushes 2, 3, 5, 9 and 17, where a push costs one write and as many copies as there are elements; every other
push costs 1. The spikes double, but so does the gap between them.}
\end{omfigure}

\begin{omfigure}
\begin{tikzpicture}
\begin{axis}[width=12cm, height=4.6cm, xmin=0.3, xmax=17.7, ymin=-2, ymax=18, xtick={1,2,3,5,9,17},
  xlabel={push number}, ylabel={$\Phi$}, tick label style={font=\small}]
\addplot[thick, omDef, mark=*, mark size=1.2pt] table[x=i, y=phi]
  {code/bachelor-1/14-dynamic-arrays/figures/copies.dat};
\addplot[only marks, mark=*, mark size=1.2pt, omDef] coordinates {(0.5,-1)};
\node[font=\scriptsize, align=left, anchor=north west] at (axis cs:1,17.5) {a cheap push: actual 1, $\Delta\Phi = +2$\\a
resize at size $s$: actual $1 + s$,\\$\Delta\Phi = 2 - s$; both sum to 3};
\end{axis}
\end{tikzpicture}
\par\smallskip{\small The potential $\Phi = 2 \cdot \text{size} - \text{capacity}$ along the same pushes: it rises by 2
at each cheap push, storing credit, and falls to 2 at each resize, spending it on the copies. The first dot is the empty \omterm{def:b1:pointers-and-arrays:array}{array}, $\Phi = -1$.}
\end{omfigure}

The same analysis applies to any \omterm{def:b1:dynamic-memory:growth}{growth factor} $r > 1$, with $r/(r - 1)$ copies per push in place of 2
(\cref{exo:b1:dynamic-arrays:10}); but not to a growth by a constant number of cells, which makes the copies quadratic.
C++'s \code{std::vector} specifies \code{push_back} as amortised constant time, which rules out a constant
increment.

\omoutput{code/bachelor-1/14-dynamic-arrays/outputs/copies.txt}{Element copies for $n$ pushes, and per push, for
growth by 1.5, by 2 and by 16 cells.}\label{out:b1:dynamic-arrays:copies}

\omoutput{code/bachelor-1/14-dynamic-arrays/outputs/timing.txt}{The same three policies timed on the author's laptop
(best of three, nanoseconds per push).}\label{out:b1:dynamic-arrays:timing}

Copies per push oscillate with $n$ (they depend on where $n$ falls between two resizes) but stay below 2 for doubling and
below 3 for the factor 1.5, as the bound says; with 16 more cells each time, they grow like $n/32$. The times agree: doubling
costs under half a nanosecond per push and the factor 1.5 from half a nanosecond to under three, while the increment
costs four to six times more at each fourfold increase of $n$, as copies growing like $n/32$ predict.

The binary counter is the second classic example. A counter of $K$ bits is incremented by turning its trailing 1s into
0s and the next 0 into a 1. An increment can flip all $K$ bits, but bit 0 flips at every increment, bit 1 at every
second, bit $j$ at every $2^j$-th: $n$ increments make fewer than $n(1 + \frac12 + \frac14 + \cdots) = 2n$ flips
(\cref{exo:b1:dynamic-arrays:8,exo:b1:dynamic-arrays:9}).

\omcode{code/bachelor-1/14-dynamic-arrays/listings/counter.c}{6}{11}{An increment, counting the flipped
bits}\label{lst:b1:dynamic-arrays:counter}

\omoutput[1-7]{code/bachelor-1/14-dynamic-arrays/outputs/counter.txt}{Total flips for $n$ increments from
zero.}\label{out:b1:dynamic-arrays:counter}

\subsection*{Shrinking}\label{sec:b1:dynamic-arrays:shrinking}

An \omterm{def:b1:pointers-and-arrays:array}{array} that only grows wastes memory after many pops. Halving the \omterm{def:b1:dynamic-arrays:dynamic-array}{capacity} when the \omterm{def:b1:pointers-and-arrays:array}{array} is half full looks
symmetric with doubling, and is a trap: at a \omterm{def:b1:dynamic-arrays:dynamic-array}{size} just above a power of 2, a push doubles the \omterm{def:b1:dynamic-arrays:dynamic-array}{capacity}, a pop halves it,
and each copies every element. Shrinking at a quarter leaves room in both directions after each resize.

\omoutput{code/bachelor-1/14-dynamic-arrays/outputs/thrash.txt}{After filling to 1024 elements, 1000 alternating pushes
and pops under each shrinking policy.}\label{out:b1:dynamic-arrays:thrash}

\section{The circular buffer}\label{sec:b1:dynamic-arrays:ring}

A queue takes elements at the front and adds them at the back. In an \omterm{def:b1:pointers-and-arrays:array}{array}, taking the front would shift every element,
or leave the used part drifting to the right; a \omterm{def:b1:dynamic-arrays:ring}{circular buffer} lets it wrap around.

\begin{definition}[Circular buffer]\label{def:b1:dynamic-arrays:ring}
A \emph{circular buffer}\index{circular buffer} stores a queue in an \omterm{def:b1:pointers-and-arrays:array}{array} of \omterm{def:b1:dynamic-arrays:dynamic-array}{capacity} $c$ with a head index $h$ and a
count $k$: the queue is $a[h], a[(h + 1) \bmod c], \dots, a[(h + k - 1) \bmod c]$. Taking the front advances $h$ modulo
$c$; adding at the back writes $a[(h + k) \bmod c]$.
\end{definition}

\omcode{code/bachelor-1/14-dynamic-arrays/listings/ring.c}{34}{51}{Adding and taking, with growth when
full}\label{lst:b1:dynamic-arrays:ring}

\omoutput{code/bachelor-1/14-dynamic-arrays/outputs/ring_demo.txt}{A queue in five cells: the front, in brackets, moves
right, the back wraps to the start, and the eighth element makes it grow.}\label{out:b1:dynamic-arrays:ring}

\begin{omfigure}
\begin{tikzpicture}[x=8mm, y=8mm]
\foreach \v/\i in {6/0, ./1, 3/2, 4/3, 5/4} {
  \draw (\i,0) rectangle (\i+1,1); \node[font=\small] at (\i+0.5,0.5) {\v};}
\fill[omDef!25] (2,0) rectangle (3,1); \node[font=\small] at (2.5,0.5) {3};
\draw[-{Stealth}, thick, omThm] (2.5,1.7) -- (2.5,1.05); \node[font=\scriptsize, above] at (2.5,1.7) {head $h = 2$};
\draw[-{Stealth}, thick, omCode] (1.1,1.7) -- (1.5,1.05); \node[font=\scriptsize, anchor=south east] at (1.3,1.7)
{next write $(h + k) \bmod 5 = 1$};
\draw[-{Stealth}, dashed] (5.05,0.5) .. controls (6.2,0.5) and (6.2,-1.6) .. (2.5,-1.6) .. controls (-1.2,-1.6) and
(-1.2,0.5) .. (-0.05,0.5);
\node[font=\scriptsize] at (2.5,-2.0) {$k = 4$: the queue is 3, 4, 5, 6};
\begin{scope}[xshift=7.2cm]
\foreach \v/\i in {3/0, 4/1, 5/2, 6/3, 7/4, 8/5, ./6, ./7, ./8, ./9} {
  \draw (\i*0.7,0) rectangle (\i*0.7+0.7,1); \node[font=\small] at (\i*0.7+0.35,0.5) {\v};}
\fill[omDef!25] (0,0) rectangle (0.7,1); \node[font=\small] at (0.35,0.5) {3};
\node[font=\scriptsize, above] at (3.5,1.1) {after growing: copied in queue order, $h = 0$};
\end{scope}
\end{tikzpicture}
\par\smallskip{\small A \omterm{def:b1:dynamic-arrays:ring}{circular buffer} before and after it grows. On the left, the queue 3, 4, 5, 6 has wrapped: it
starts at index 2 and continues at index 0. Growing copies the queue in order to the start of a larger \omterm{def:b1:pointers-and-arrays:array}{array}.}
\end{omfigure}

Keeping the count rather than a second index matters: with a head and a tail index only, the full and the empty
buffer both have $h = t$ (\cref{exo:b1:dynamic-arrays:5}). Growth copies the elements in queue order, from the head,
not in the order of the cells: copying the cells as they are would break the wrap-around.

\begin{inpractice}
Python's \code{list} and C++'s \code{std::vector} are \omterm{def:b1:dynamic-arrays:dynamic-array}{dynamic arrays}. Python's \code{collections.deque} is a
chain of \omterm{def:b1:dynamic-arrays:dynamic-array}{fixed-size} blocks, and the buffers between a producer and a consumer, from keyboard input to network cards,
are queues in fixed blocks with wrapping indices. A program that knows how many elements it will push can reserve the \omterm{def:b1:dynamic-arrays:dynamic-array}{capacity} once and avoid every
copy.
\end{inpractice}

\section{Exercises}\label{sec:b1:dynamic-arrays:exercises}

\begin{exercise}[$\star$]\label{exo:b1:dynamic-arrays:1}
Predict the three lines this program prints, with the doubling \omterm{def:b1:pointers-and-arrays:array}{array} of \cref{lst:b1:dynamic-arrays:vec}.
\omcode{code/bachelor-1/14-dynamic-arrays/listings/exo_ops.c}{6}{18}{Pushes and pops}
\end{exercise}

\begin{exercise}[$\star$]\label{exo:b1:dynamic-arrays:2}
How many element copies do 10 pushes make on the doubling \omterm{def:b1:pointers-and-arrays:array}{array}, starting from \omterm{def:b1:dynamic-arrays:dynamic-array}{capacity} 1? And 10 pushes when each
resize adds 2 cells?
\end{exercise}

\begin{exercise}[$\star$]\label{exo:b1:dynamic-arrays:3}
A \omterm{def:b1:dynamic-arrays:ring}{circular buffer} of \omterm{def:b1:dynamic-arrays:dynamic-array}{capacity} 4 starts empty. Draw its cells, head and count after: put 1, put 2, put 3, take, take,
put 4, put 5, put 6, take.
\end{exercise}

\begin{exercise}[$\star$]\label{exo:b1:dynamic-arrays:4}
How many bits do 8 increments flip, starting from 0? List the flips of each increment.
\end{exercise}

\begin{exercise}[$\star\star$]\label{exo:b1:dynamic-arrays:5}
This buffer keeps a head and a tail index. Explain the output, and give two ways to repair it.
\omcode{code/bachelor-1/14-dynamic-arrays/listings/exo_ring_bug.c}{4}{19}{Head and tail only}
\omoutput[2-3]{code/bachelor-1/14-dynamic-arrays/outputs/exo_ring_bug.txt}{Four puts, and nothing to take.}
\end{exercise}

\begin{exercise}[$\star\star$]\label{exo:b1:dynamic-arrays:6}
In \cref{out:b1:dynamic-arrays:thrash}, explain the 1024 copies per operation of the half policy, and the one copy per
operation of the quarter policy. Give a sequence of $m$ operations that costs $\Theta(m^2)$ copies under the half
policy.
\end{exercise}

\begin{exercise}[$\star\star$]\label{exo:b1:dynamic-arrays:7}
Write an abstract type ``stack of int'' (\code{stack.h}, \code{stack.c}) on top of \code{vec.h}, with a \omterm{def:b1:dynamic-arrays:operations}{constructor}, two
\omterm{def:b1:dynamic-arrays:operations}{accessors} (\code{stack_empty}, \code{stack_top}) and two \omterm{def:b1:dynamic-arrays:operations}{transformers}, and test it, including on a bracket checker.
\end{exercise}

\begin{exercise}[$\star\star$]\label{exo:b1:dynamic-arrays:8}
By the \omterm{def:b1:dynamic-arrays:methods}{aggregate method}, prove that $n$ increments of a binary counter, from 0, flip fewer than $2n$ bits.
\end{exercise}

\begin{exercise}[$\star\star$]\label{exo:b1:dynamic-arrays:9}
Prove the same bound by the \omterm{def:b1:dynamic-arrays:methods}{accounting method}: how many credits per increment, and where are they stored?
\end{exercise}

\begin{exercise}[$\star\star\star$]\label{exo:b1:dynamic-arrays:10}
The \omterm{def:b1:pointers-and-arrays:array}{array} grows by a factor $r > 1$ when full (\omterm{def:b1:dynamic-arrays:dynamic-array}{capacity} $c$ becomes at least $rc$, for
instance $\lceil rc \rceil$; the factor 1.5 of the programs gives $c + \lfloor c/2 \rfloor + 1$). Prove with a potential that $n$
pushes make at most about $\frac{r}{r - 1} n$ copies, and compare with \cref{out:b1:dynamic-arrays:copies}.
\end{exercise}

\begin{exercise}[$\star\star\star$]\label{exo:b1:dynamic-arrays:11}
With growth $\times 2$ and shrinking to half when a quarter full, prove that any sequence of $m$ pushes and pops from the
empty \omterm{def:b1:pointers-and-arrays:array}{array} makes at most $2m$ copies. Use $\Phi = 2s - c$ when $s \geq c/2$ and $\Phi = c/2 - s$ when $s < c/2$, for \omterm{def:b1:dynamic-arrays:dynamic-array}{size}
$s$ and \omterm{def:b1:dynamic-arrays:dynamic-array}{capacity} $c$.
\end{exercise}

\begin{exercise}[$\star\star\star$]\label{exo:b1:dynamic-arrays:12}
Write a \omterm{def:b1:dynamic-arrays:dynamic-array}{dynamic array} of elements of any type: the \omterm{def:b1:dynamic-arrays:operations}{constructor} takes the size of an element in bytes, push copies an
element from an address, and \code{gvec_at} returns the address of an element. Test it with \code{double} and with a
structure.
\end{exercise}

\section{Problem: The binary counter and the elastic array}\label{sec:b1:dynamic-arrays:problem}

\begin{problem}[{Weekend problem --- The binary counter and the elastic array}]\label{pb:b1:dynamic-arrays:1}
A program tries every sequence of $m$ pushes and pops on the policy of \code{vec.c} (double when full, halve when a
quarter full), for even $m$ up to 26, and prints the largest number of copies with one sequence that makes them.
\omcode{code/bachelor-1/14-dynamic-arrays/listings/elastic.c}{5}{24}{Copies of one sequence of
operations}\label{lst:b1:dynamic-arrays:elastic}
\omoutput{code/bachelor-1/14-dynamic-arrays/outputs/elastic.txt}{The worst sequences found.}\label{out:b1:dynamic-arrays:elastic}
\omoutput[8-8]{code/bachelor-1/14-dynamic-arrays/outputs/counter.txt}{The counter of \cref{lst:b1:dynamic-arrays:counter},
with a decrement, alternating around $2^{15}$.}\label{out:b1:dynamic-arrays:counter-dec}

\needspace{4\baselineskip}
\medskip\noindent\textbf{Part I --- The counter.}
\begin{enumerate}
\item Show that an increment from $x$ flips $t + 1$ bits, where $t$ is the number of trailing 1s of $x$ (when $t < K$).
\item How many of the integers $0 \leq x < 2^K$ have exactly $t$ trailing 1s?
\item Show that bit $j$ flips exactly $\lfloor n/2^j \rfloor$ times in $n$ increments from 0.
\item Deduce that the total is $\sum_j \lfloor n/2^j \rfloor < 2n$, and that it equals $2n - 1$ when $n$ is a power of 2.
\item Check the table of \cref{out:b1:dynamic-arrays:counter}.
\item What is the worst case of a single increment, and why does it not contradict the amortised bound?
\end{enumerate}

\needspace{4\baselineskip}
\medskip\noindent\textbf{Part II --- The doubling array, three ways.}
\begin{enumerate}[resume]
\item Count exactly the copies made by $n$ pushes from \omterm{def:b1:dynamic-arrays:dynamic-array}{capacity} 1, as a function of $n$.
\item Write the accounting argument for pushes with 3 credits each, and say why 2 credits would not do.
\item Check that $\Phi = 2 \cdot \text{size} - \text{capacity}$ is at least 0 after every push.
\item Give the \omterm{def:b1:dynamic-arrays:amortised}{amortised cost} of each push of \cref{out:b1:dynamic-arrays:demo} with this $\Phi$.
\item With $n$ pushes, what is the largest \omterm{def:b1:dynamic-arrays:dynamic-array}{capacity}, and so the largest unused memory, as a fraction of the \omterm{def:b1:dynamic-arrays:dynamic-array}{size}?
\item Which of the three methods would you use to explain the result to a colleague, and why?
\end{enumerate}

\needspace{4\baselineskip}
\medskip\noindent\textbf{Part III --- Shrinking.}
\begin{enumerate}[resume]
\item With shrinking at a half, describe the state after filling to $2^k$ elements, then one push, then one pop.
\item Deduce that $m$ alternating operations after the fill cost $\Theta(m \cdot 2^k)$ copies.
\item With shrinking at a quarter, show that after any resize the \omterm{def:b1:dynamic-arrays:dynamic-array}{size} is exactly half the new \omterm{def:b1:dynamic-arrays:dynamic-array}{capacity}.
\item Deduce that between two resizes at least a quarter of the \omterm{def:b1:dynamic-arrays:dynamic-array}{capacity}'s worth of operations happen.
\item Read the table of \cref{out:b1:dynamic-arrays:elastic}: which sequences are the worst, and what is the largest
ratio of copies to operations?
\item Show that $2^k + 1$ pushes make $2^{k+1} - 1$ copies, so that the ratio approaches 2.
\end{enumerate}

\needspace{4\baselineskip}
\medskip\noindent\textbf{Part IV --- A decrement breaks the bound.}
\begin{enumerate}[resume]
\item Write the decrement of the counter, and give the flips of a decrement from $2^j$.
\item From $2^{15} - 1$, alternate increment and decrement: how many flips per operation? Check with
\cref{out:b1:dynamic-arrays:counter-dec}.
\item Deduce that $n$ operations among increments and decrements can cost $\Theta(n \log n)$ flips when the counter has
$\log_2 n$ bits.
\item With digits $-1$, 0 and 1 in base 2 (each number then has several representations), an increment and a decrement
can each be made to change a bounded number of digits in the amortised sense. Why does this not contradict the previous
question?
\item Is the amortised bound of the \omterm{def:b1:pointers-and-arrays:array}{array}, $2m$ copies, still valid for any mix of pushes and pops? Why do pops not break
it as decrements break the counter's?
\item State the named result of the problem, \emph{the elastic \omterm{def:b1:pointers-and-arrays:array}{array} theorem}: the constant $c$ with at most $c \cdot m$
copies for any $m$ operations, why it is the smallest, and the worst ratios found by the program.
\end{enumerate}
\end{problem}

